\section{Conclusion}
\label{sec:conclusion}

This paper asks a simple but previously difficult question: what does winning the H-1B lottery do to the individual worker? I answer that question by combining recent deidentified H-1B administrative records with LinkedIn-based worker histories from Revelio Labs and by developing a probabilistic linkage strategy that can connect the two without direct personal identifiers. The approach is designed around the institutional reality that many recent H-1B applicants are already working for the sponsoring employer before the lottery, often under OPT. That fact makes the worker-level question empirically tractable.

The paper makes two main contributions. First, it introduces a linkage framework for deidentified immigration records that should be useful beyond this specific application. The contribution is not only that a link can be made, but that the uncertainty in the link can be carried transparently into the analysis through match weights and high-precision subsamples. Second, it provides the first worker-level causal evidence, to my knowledge, on the consequences of winning the H-1B lottery. The current results indicate that lottery winners are more likely to remain with the sponsoring employer and more likely to remain in the United States, but the magnitudes are small.

Those small effects are themselves informative. They suggest that, for a large subset of recent applicants, the H-1B cap is not simply a gate that cleanly separates entry from non-entry. Instead, it interacts with a broader system of temporary work authorization in which many lottery losers continue to work, reapply, extend OPT, or transition through other channels. For the modern pre-H-1B worker, winning appears to increase labor-market stability more than it changes the extensive margin of being present in the United States at all.

This interpretation has two policy implications. The first is that reforms to the H-1B system should be evaluated not only in terms of how many workers they admit, but also in terms of how they change employer dependence, repeated application costs, and uncertainty for workers already trained and employed in the United States. The second is that the individual-level incidence of H-1B rationing may differ from the incidence suggested by firm-level studies: the main worker-level effects may operate through retention and mobility rather than through abrupt displacement.

There are several natural extensions. One is to study wage and career progression in more detail once the preferred compensation measures are finalized. Another is to connect the linked applicant panel to productivity outcomes, such as GitHub activity or entrepreneurial transitions. A third is to study return migration and third-country reallocation more directly. More broadly, the project shows that the deidentified nature of administrative immigration data, while a serious obstacle, need not prevent credible worker-level analysis when combined with modern linkage methods and a careful understanding of institutional context.
